\subsection{Análisis de resultados y conclusiones}
\label{subsec:analisis-de-resultados-y-conclusiones}

\subsubsection{Práctica 1 --- Ley de Ohm}

\todo{P1 --- Análisis (2.5.1--2.5.4): compatibilidad del ajuste con $V=IR$, ¿es $b$ compatible con cero?, ¿es preferible la regresión a promediar $R_n=V_n/I_n$?; consecuencias de $b \neq 0$; hipótesis (contacto, temperatura, no linealidad); mejoras.}

% ####################################################################################################################
% ####################################################################################################################

\subsubsection{Práctica 2 --- Comparación entre Modelo y Sistema}

Los tres valores medidos son compatibles con el valor nominal dentro de la tolerancia de $\pm 5\%$ declarada por el fabricante.
La Ley de Ohm establece $V = IR$, que sustenta el principio del ohmímetro; los resultados son consistentes con ese modelo.

\subsubsection{Práctica 2 --- Consecuencia de las Discrepancias}

Las pequeñas diferencias entre el valor nominal y el medido (por ejemplo, R3: 5{,}50 k$\Omega$ frente a nominal 5{,}6 k$\Omega$) son normales y esperadas: la tolerancia del $\pm 5\%$ indica que la resistencia puede fabricarse con ese margen de variación.
No representan un error experimental sino variación en el proceso de fabricación.

\subsubsection{Práctica 2 --- Hipótesis de las Discrepancias}

Las fuentes de incerteza en la medición son:
\begin{itemize}
    \item \textbf{Error del instrumento:} cuantificado por la fórmula $\Delta R = \%rdg \times R + dgt \times resolución$.
    \item \textbf{Variación de fabricación:} las resistencias se fabrican dentro de la banda de tolerancia declarada.
    \item \textbf{Temperatura:} la resistencia de un conductor varía con la temperatura; a temperatura ambiente el efecto es pequeño pero no nulo.
\end{itemize}

\subsubsection{Práctica 2 --- Mejoras al Método}

\begin{itemize}
    \item Usar siempre el \textbf{rango más ajustado} que contenga la medición para minimizar la incerteza (como se hizo con R3, pasando del rango 2000 k$\Omega$ al 20 k$\Omega$).
    \item Asegurarse de que las puntas no toquen más de una resistencia a la vez para evitar medir combinaciones serie/paralelo no deseadas.
\end{itemize}

Se midieron exitosamente tres resistencias con el multímetro digital, obteniendo resultados compatibles con los valores nominales dentro de la tolerancia declarada.
La práctica ilustró la importancia de seleccionar el rango adecuado del instrumento para minimizar la incerteza, y de estimar correctamente dicha incerteza a partir de las especificaciones del fabricante.

% ####################################################################################################################
% ####################################################################################################################

\subsubsection{Práctica 3 --- Influencia del aparato de medida}

\todo{P3 --- Análisis (2.5.1--2.5.4): ¿se cumple Kirchhoff?, ¿para qué valores de $R$ sí y para cuáles no?; efecto de $R_{int}=10\ \text{M}\Omega$ sobre el circuito; modelo $R'_2 = R_2 \| R_{int}$; mejoras.}
